\documentclass[11pt]{article}
\usepackage[margin=2.5cm]{geometry}
\usepackage{booktabs}
\usepackage{array}
\usepackage{hyperref}

\title{FinLingo++ -- Results}
\author{Vedika Thirumalai Nambi/24222437}
\date{17 July 2026}

\begin{document}
\maketitle

\section{Pipeline}
\begin{table}[h!]
\centering
\begin{tabular}{clll}
\toprule
Stage & What it does & Model used \\
\midrule
S1 & Document parsing & PyMuPDF + OCR \\
S2 & Simplify clause & Qwen2.5-1.5B + QLoRA \\
S3 & Retrieve evidence & FAISS + Ollama embeddings \\
S4 & Rerank evidence & MiniLM-L-12 cross-encoder \\
S5 & Classify risk & Qwen2.5-1.5B, prompted, 6 risk classes \\
S6 & Check faithfulness & DeBERTa-v3 NLI + QLoRA \\
S7 & Generate report & FastAPI + React \\
\bottomrule
\end{tabular}
\end{table}
\noindent\textit{Note: S5 uses the same Qwen2.5-1.5B model as S2 (its own \texttt{risk\_model}
setting is unset, so it falls back to the S2 simplifier model) rather than a separate larger
model. S5 is prompted only, with no fine-tuning, by explicit design.}

\section{Status checklist}
\begin{table}[h!]
\centering
\begin{tabular}{ll}
\toprule
Item & Status \\
\midrule
S2 domain warm-up training & Done \\
S2 simplifier adapter training & Done \\
S6 DeBERTa verifier training (v3, 5,000 examples) & Done \\
Verifier calibration \& FaithBench reproduction (n=500) & Done \\
FinanceBench retrieval-query-mode ablation & Done \\
Reranker ablation (L-6 vs L-12) & Done \\
Multi-verifier comparison (DeBERTa v3 vs ModernBERT) & Done \\
Failure analysis & Done \\
Test suite (111 tests) \& frontend build & Passing \\
6-class FLB benchmark (achievable scale, 91 clauses / 182 rows) & Done (built and locked; not the originally-scoped 500-clause version) \\
Full ablation suite (6 pipeline variants + 5 baselines) & Done \\
\bottomrule
\end{tabular}
\end{table}

\section{Results}

\subsection{S6 verifier: base vs.\ fine-tuned (FaithBench, 500 test examples)}
\begin{table}[h!]
\centering
\begin{tabular}{lrr}
\toprule
Metric & Base DeBERTa & Fine-tuned v3 \\
\midrule
Accuracy & 0.470 & 0.508 \\
Balanced accuracy & 0.518 & 0.543 \\
Unsupported (hallucination) recall & 0.086 & 0.234 \\
Supported F1 & 0.614 & 0.606 \\
\bottomrule
\end{tabular}
\end{table}
\noindent Fine-tuning improves accuracy, balanced accuracy, and unsupported-claim recall (nearly
tripling it). The improvement in balanced accuracy is statistically significant: paired bootstrap
$+0.038$, 95\% CI $[0.006, 0.072]$, $p=0.02$ ($n=500$). Both configurations still miss the
majority of unsupported claims (recall well under 0.5).

\subsection{S2 simplifier: with adapter vs.\ zero-shot base model (same 91-clause benchmark)}
\begin{table}[h!]
\centering
\begin{tabular}{lrr}
\toprule
Metric & Zero-shot base Qwen & With adapter (deployed) \\
\midrule
SARI & 49.01 & 49.44 \\
BERTScore F1 & 0.895 & 0.895 \\
Mean FK grade & 9.68 & 9.77 \\
FK grade $\leq 9$ rate & 44.0\% & 41.8\% \\
\bottomrule
\end{tabular}
\end{table}
\noindent The adapter produces a small, real improvement in SARI with essentially no change in
BERTScore, and does not clearly improve FK grade over the zero-shot base model on this benchmark
-- both configurations fail the FK $\leq 9$ target by a similar margin. This is measured on the
full 91-clause official benchmark (larger and more representative than a small held-out dev
split), using the metrics tracked by the current evaluation harness (SARI/BERTScore/FK); a
semantic-entailment/acceptance-rate breakdown is not part of the current evaluator's output.

\subsection{Retrieval and reranking (FLB, 91 queries)}
\begin{table}[h!]
\centering
\begin{tabular}{lrrr}
\toprule
System & Recall@5 & MRR & Latency \\
\midrule
FAISS only (pre-rerank, Recall@8) & 0.111 & 0.029 & -- \\
+ MiniLM-L-6 & 0.067 & 0.021 & 46 ms \\
+ MiniLM-L-12 (deployed) & 0.067 & 0.031 & 81 ms \\
\bottomrule
\end{tabular}
\end{table}
\noindent Reranking does not improve recall@5 over pre-rerank recall@8 on this benchmark (different
$k$, so not directly comparable), but L-12 improves ranking quality (MRR) over L-6 at roughly
1.8$\times$ the latency. Absolute retrieval numbers on this benchmark are far lower than on
FinanceBench (see \S3.3 of the earlier baseline) -- confirmed to be a corpus/query-match
bottleneck specific to this benchmark, not a retrieval-architecture defect (see \S3.3 below).

\subsection{Query type used for retrieval (FinanceBench, 30 examples)}
\begin{table}[h!]
\centering
\begin{tabular}{lrr}
\toprule
Query & Recall@8 & Recall@5 (post-rerank) \\
\midrule
Original clause & 0.883 & 0.700 \\
Simplified clause & 0.772 & 0.656 \\
Combined & 0.867 & 0.706 \\
\bottomrule
\end{tabular}
\end{table}
\noindent Simplifying the clause before retrieval measurably lowers recall relative to using the
original clause -- confirms simplification strips retrieval-relevant terms. Recall on this
FinanceBench benchmark (0.70--0.88) is far higher than recall on the FLB benchmark (0.07--0.11,
\S3.3) using the identical embedding and reranker models: this cross-corpus gap is the confirmed
root cause of FLB's own low retrieval numbers, not a defect in the retrieval architecture itself.

\section{Main findings}
\begin{itemize}
\item Reranking improves ranking quality (MRR) but not recall@k on the FLB benchmark; on
  FinanceBench, reranking and the choice of query text both measurably affect results.
\item Simplifying text before retrieval loses key financial/regulatory terms and hurts
  retrieval -- confirmed on FinanceBench and identified as the dominant driver of FLB's own
  faithfulness-verification bottleneck.
\item The S2 adapter gives a small, real SARI improvement over the zero-shot base model, with no
  clear FK-grade benefit on the full 91-clause benchmark.
\item The S6 verifier improves significantly with fine-tuning (statistically significant on
  balanced accuracy, unsupported-recall nearly tripled) but both configurations still miss the
  majority of unsupported claims.
\item The primary (chunk-grounded) faithfulness metric is near-zero on FLB. Root cause: the
  chunk judged to actually ground a hypothesis is present in the verifier's evidence set only
  6.6\% of the time (12/182) -- a retrieval-coverage bottleneck, not a benchmark-labelling
  defect or a verifier failure in isolation. This is disclosed, not smoothed over.
\item The risk classifier (S5) macro-F1 (0.56) remains below the 0.75 target. This gap reflects
  a genuine tension in the project's own design: S5 is required to be prompted-only, no
  fine-tuning, while targeting a fine-tuning-level score. Training a classifier to close this
  gap was considered and declined, since it would depart from the specified method.
\item The full-scale 91-clause, 6-class FLB benchmark was successfully built and locked this
  round (up from the earlier 3-class-only blocked state); it remains smaller than the originally
  planned 500-clause version due to genuine raw-data and quality-gate supply limits, not a
  relaxed target -- documented per-class in the benchmark plan.
\item The full ablation suite (isolating each pipeline component's contribution, plus five
  comparator baselines: rule-based, plain-RAG, frozen FinBERT, RAGAS-style paragraph judge, and
  SelfCheckGPT-style self-consistency) is now complete. The deployed verifier's secondary
  faithfulness metric (precision 0.906, recall 0.637) beats both alternative faithfulness-checking
  baselines (RAGAS-style: 0.689/0.560; SelfCheckGPT-style: 0.849/0.495) on the same benchmark --
  a genuine, positive comparative finding. A frozen FinBERT baseline with a simple logistic-
  regression head reaches macro-F1 0.69 on risk classification, ahead of the deployed prompted
  classifier (0.56); verified via the codebase's existing deterministic document-level train/test
  partition that this is not a data-leakage artifact, reinforcing that S5's ``no fine-tuning''
  constraint (\S4) has a real, measurable cost.
\end{itemize}

\end{document}
